% !TEX program = pdflatex
\documentclass[11pt,a4paper]{article}
\usepackage[margin=1.02in]{geometry}
\usepackage{amsmath,amssymb,amsthm,mathtools}
\usepackage{booktabs,microtype}
\usepackage{xcolor}
\usepackage[numbers,sort&compress]{natbib}
\usepackage[colorlinks=true,linkcolor=blue!55!black,citecolor=blue!55!black,urlcolor=blue!55!black]{hyperref}

\newtheorem{theorem}{Theorem}[section]
\newtheorem{proposition}[theorem]{Proposition}
\newtheorem{lemma}[theorem]{Lemma}
\newtheorem{corollary}[theorem]{Corollary}
\theoremstyle{definition}
\newtheorem{definition}[theorem]{Definition}
\theoremstyle{remark}
\newtheorem{remark}[theorem]{Remark}
\newcommand{\R}{\mathbb{R}}
\newcommand{\Gr}{\operatorname{Gr}}
\newcommand{\sgn}{\operatorname{sgn}}
\newcommand{\diag}{\operatorname{diag}}
\newcommand{\abs}[1]{\lvert#1\rvert}
\newcommand{\norm}[1]{\lVert#1\rVert}

\title{\textbf{Generic Sign Rigidity of Real Projection Kernels from Two-Point Data}}
\author{}
\date{Revised working draft -- 29 October 2026}
\hypersetup{pdfauthor={}}

\begin{document}
\maketitle

\begin{abstract}
Let $K$ be a real symmetric orthogonal projection. Its diagonal and the moduli of its
off-diagonal entries are exactly the data contained in its principal minors of orders one
and two, but they do not display the signs of $K_{ij}$.  We prove that this apparent loss is
generically repaired by idempotency.  For every $N\ge5$ and $2\le r\le N-2$, almost every
rank-$r$ projection is determined, up to conjugation by a diagonal sign matrix, by
$\{K_{ii},|K_{ij}|\}$.  The proof separates $K^2=K$ into a diagonal modulus identity and a
linear system for gauge-invariant triangle holonomies.  A zero-subset-sum condition on the
signed wedge weights is a finite certificate of uniqueness.  We show that every forbidden
subset sum is a nonzero real-analytic function on $\Gr(r,N)$; its zero locus has measure
zero, and a finite union argument gives the theorem.  We also give a certified reconstruction
procedure: solve the finite holonomy system, test the certificate, fix a spanning-tree gauge,
and recover all signs.  The result is a theorem about finite real projection kernels.  Its
use for determinantal boundary data is only a dictionary: in a finite real pure
determinantal completion, one- and two-boundary data generically fix the higher principal
minors.  No continuum, Hermitian, or quantum-gravity completeness claim is made.
\end{abstract}

\section{Introduction}

For a real determinantal point process, inclusion probabilities are principal minors of a
symmetric contraction $K$ \citep{HKPV2006,KuleszaTaskar2012}.  The order-one and order-two
minors give
\begin{equation}
 d_i=K_{ii},\qquad \det K_{\{i,j\}}=d_id_j-K_{ij}^2,
 \label{eq:minors}
\end{equation}
and hence determine $m_{ij}=|K_{ij}|$.  They do not determine the signs.  General
principal-minor recovery therefore uses cycle data, normally supplied by minors of order at
least three \citep{UrschelEtAl2017,Brunel2018,BrunelUrschel2024}.  All principal minors are
invariant under
\begin{equation}
 K\longmapsto D K D,\qquad D=\diag(\epsilon_1,\ldots,\epsilon_N),\quad
 \epsilon_i\in\{\pm1\},
 \label{eq:gauge}
\end{equation}
so recovery can never be sharper than this diagonal-sign gauge.

We add one hypothesis: $K$ is an orthogonal projection.  Projection kernels are basic
extreme determinantal kernels \citep{Lyons2003,HKPV2006}; algebraically, the added input is
simply $K^2=K$.  The question is whether this quadratic identity replaces missing cycle
minors.  Our answer is affirmative on a full-measure subset of every interior real
Grassmannian covered below.

\begin{theorem}[Generic projection sign rigidity]\label{thm:main}
Fix $N\ge5$ and $2\le r\le N-2$.  For almost every $K\in\Gr(r,N)$, if $K'$ is a real
symmetric projection satisfying
\[
 K'_{ii}=K_{ii},\qquad |K'_{ij}|=|K_{ij}|\quad(i\ne j),
\]
then $K'=DKD$ for a diagonal sign matrix $D$.
\end{theorem}

Here ``almost every'' refers to the natural smooth measure class on $\Gr(r,N)$.  The
theorem includes generic full off-diagonal support; it is not a claim for every projection.
The uniform statement is deliberately scoped to $N\ge5$ and $2\le r\le N-2$.  Sparse
supports, exceptional subset-sum loci, complex Hermitian projections, and continuum kernels
require separate analyses.

The key observation is that the off-diagonal part of idempotency becomes \emph{linear} after
passing from edge signs to triangle holonomies.  Sign rigidity is then a finite lattice
intersection problem.  This differs from the established principal-minor results: those
recover cycle signs from higher minors, whereas here only orders one and two are given and
idempotency supplies the extra equations.

\section{Gauge-invariant form of idempotency}

Let $K=K^\top=K^2\in\R^{N\times N}$ have rank $r$.  Put
\begin{equation}
 d_i=K_{ii},\qquad m_{ij}=|K_{ij}|,
 \qquad \sigma_{ij}=\sgn K_{ij}\quad(m_{ij}>0).
\end{equation}
For the main theorem the support is generically complete.  On a complete support graph,
define the triangle holonomy
\begin{equation}
 h_{ikj}=\sigma_{ik}\sigma_{kj}\sigma_{ji}\in\{\pm1\}.
 \label{eq:h}
\end{equation}
It is invariant under \eqref{eq:gauge}.  Conversely, a valid holonomy field determines the
edge signs up to gauge: set the signs on a spanning tree, then use triangles to obtain every
remaining edge.  Valid fields obey the usual cocycle relations; not every arbitrary
$\{\pm1\}$ assignment to triangles is valid.

\begin{proposition}[Idempotency split]\label{prop:split}
For full off-diagonal support, $K^2=K$ is equivalent to
\begin{align}
 \sum_{k\ne i}m_{ik}^2&=d_i(1-d_i), && i=1,\ldots,N, \label{eq:E2}\\
 \sum_{k\ne i,j}m_{ik}m_{kj}h_{ikj}
 &=m_{ij}(1-d_i-d_j), && 1\le i<j\le N. \label{eq:E1}
\end{align}
The first system contains only two-point data; the second is linear in $h$.
\end{proposition}

\begin{proof}
The diagonal equation $K_{ii}=\sum_kK_{ik}^2$ gives \eqref{eq:E2}.  For $i\ne j$,
separating the $k=i,j$ terms in $K_{ij}=\sum_kK_{ik}K_{kj}$ gives
\[
 K_{ij}(1-d_i-d_j)=\sum_{k\ne i,j}K_{ik}K_{kj}.
\]
Divide by $\sigma_{ij}$ and use $\sigma_{ij}^{-1}=\sigma_{ji}$ to obtain
\eqref{eq:E1}.  Reversing these steps recovers all entries of $K^2=K$.
\end{proof}

Write \eqref{eq:E1} as $Ah=R$, where
\begin{equation}
 A_{(ij),(ikj)}=m_{ik}m_{kj}\ge0,
 \qquad R_{ij}=m_{ij}(1-d_i-d_j).
 \label{eq:linear}
\end{equation}
Let $\mathcal H$ denote the finite set of valid triangle-holonomy fields and set
$\mathcal H_R=\{h\in\mathcal H:Ah=R\}$.  The existence assumption makes
$\mathcal H_R$ nonempty.  Since valid triangle holonomies determine a complete-support
edge-sign field up to diagonal-sign gauge, uniqueness for the given data is exactly
$|\mathcal H_R|=1$.  Equivalently, for any fixed $h_0\in\mathcal H_R$,
\begin{equation}
 \ker A\cap\{h-h_0:h\in\mathcal H\}=\{0\}.
 \label{eq:lattice}
\end{equation}
The stronger condition $\ker A\cap(\mathcal H-\mathcal H)=\{0\}$ asserts global
injectivity on all valid holonomy fields and is sufficient for uniqueness at every
feasible right-hand side.

\begin{remark}[A unique observed fiber does not imply uniform injectivity]
The distinction occurs for actual full-support projection data.
Take vertices $0,\ldots,5$ and $K=\mathbf 1\mathbf 1^\top/6$.
Then $K^2=K$, its rank is one, and $d_i=m_{ij}=1/6$.
Every target component is $R_{ij}=4/36$.
Each row of $Ah$ is a sum of four numbers $\pm1/36$, so feasibility at this $R$
forces every triangle holonomy to be $+1$; thus $\mathcal H_R$ is a singleton.

Nevertheless $A$ is not injective on all of $\mathcal H$.
Fix $\sigma_{0i}=+1$ and, among the other ten edges, let the positive edges be
\[
 S_1^+=\{13,15,23,24\},\qquad S_2^+=\{13,14,23,25\},
\]
with all remaining free-edge signs negative; $ij$ denotes the edge $\{i,j\}$.
The resulting valid triangle fields $h^{(1)},h^{(2)}$ are distinct but have identical
edge-row sums. In lexicographic order on all fifteen edges,
\[
36Ah^{(1)}=36Ah^{(2)}
 =\big(0,0,0,-2,-2,0,0,2,2,0,2,2,-2,-2,0\big)^\top.
\]
This common right-hand side is not $R$, so the collision gives no rival projection
for the observed data.
The example concerns the ancillary full-support characterization at rank one,
outside the interior-rank range of Theorem~\ref{thm:main}; that generic theorem is unchanged.
A portable exact check enumerates all $2^{10}$ star-gauge sign assignments and confirms
both the singleton target fiber and this non-target collision.
\end{remark}

\section{A finite uniqueness certificate}

For a fixed projection, define the signed wedge weights in row $(i,j)$ by
\begin{equation}
 w_k^{(ij)}=m_{ik}m_{kj}h_{ikj}=K_{ik}K_{kj}\sigma_{ij},\qquad k\ne i,j.
 \label{eq:w}
\end{equation}

\begin{lemma}[Zero-subset-sum certificate]\label{lem:zss}
Suppose $K$ has full off-diagonal support and, for every edge $(i,j)$, no nonempty subset
of $\{w_k^{(ij)}:k\ne i,j\}$ has sum zero.  Then every projection with the same diagonal
and entrywise moduli is gauge-equivalent to $K$.
\end{lemma}

\begin{proof}
Compare $K$ with a rival sign field.  Let $S$ be the set of flipped edges.  A triangle is
flipped exactly when it contains an odd number of edges of $S$.  Subtracting the two copies
of \eqref{eq:E1} gives, in each row,
\[
 \sum_{k:\,(i,k,j)\text{ is flipped}}w_k^{(ij)}=0.
\]
The hypothesis makes every such subset empty, so no triangle is flipped.  In the complete
graph, triangles span the cycle space over $\mathbb F_2$.  An edge set meeting every
triangle evenly is therefore in the cut space.  Hence $S$ is a vertex cut, precisely the
gauge transformation \eqref{eq:gauge}.
\end{proof}

This certificate is sufficient, not necessary.  It can be checked exactly using
meet-in-the-middle subset sums in $O(N^2 2^N)$ arithmetic operations.  It is a
\emph{post-recovery} certificate: the weights use a candidate holonomy field.  The data-only
problem remains finite because one can solve \eqref{eq:linear} over valid cocycles (or
enumerate a gauge-fixed edge-sign field), and then apply the certificate to each survivor.

\section{Proof of generic rigidity}

We now show that the certificate holds almost everywhere.  For an edge $(i,j)$ and a
nonempty $T\subseteq[N]\setminus\{i,j\}$, define
\begin{equation}
 F_{ij,T}(K)=\sum_{k\in T}K_{ik}K_{kj}.
 \label{eq:F}
\end{equation}
Up to the common nonzero factor $\sigma_{ij}$, this is exactly the corresponding subset sum
of \eqref{eq:w}.

\begin{lemma}[Nonvanishing]\label{lem:nv}
For $N\ge5$ and $2\le r\le N-2$, no function $F_{ij,T}$ is identically zero on
$\Gr(r,N)$.
\end{lemma}

\begin{proof}
Choose $k\in T$ and set $u=(e_i+e_j+e_k)/\sqrt3$.  Since $r-1\le N-3$, choose distinct
indices $a_1,\ldots,a_{r-1}$ outside $\{i,j,k\}$.  Let $K$ be the orthogonal projection
onto
\[
 \operatorname{span}\{u,e_{a_1},\ldots,e_{a_{r-1}}\}.
\]
Then $K_{ik}=K_{kj}=1/3$, while $K_{i\ell}K_{\ell j}=0$ for every $\ell\ne k$.
Consequently $F_{ij,T}(K)=1/9$.  The witness need not have full support; it only proves
that the analytic function is nonzero.
\end{proof}

\begin{proof}[Proof of Theorem~\ref{thm:main}]
The entries $K_{ab}$ and the finitely many functions $F_{ij,T}$ are real analytic on the
connected analytic manifold $\Gr(r,N)$.  Each entry is not identically zero, so the union
of all zero-entry loci has measure zero; outside it, $K$ has full support.  By
Lemma~\ref{lem:nv}, every $F_{ij,T}$ is nonzero as an analytic function.  The zero set of a
nonzero real-analytic function has measure zero.  There are finitely many triples
$(i,j,T)$, hence almost every $K$ is simultaneously full-support and free of zero subset
sums.  Lemma~\ref{lem:zss} then gives uniqueness up to gauge.
\end{proof}

The proof is qualitative but constructive enough to yield exact rational witnesses.  If
$V\in\mathbb Q^{r\times N}$ has full row rank, then
$K=V^\top(VV^\top)^{-1}V$ is rational.  Clearing denominators turns every certificate row
into an integer subset-sum test.  Such checks verify instances; the analytic argument above,
not the finite sample, proves the almost-everywhere statement.

\section{Certified reconstruction}

Given exact $(d,m)$ and the information that a real rank-$r$ projection exists, a complete
finite procedure is:
\begin{enumerate}
\item Check $\sum_i d_i=r$ and the modulus equations \eqref{eq:E2}.
\item Fix a spanning-tree gauge (for complete support, take $\sigma_{1i}=+1$).
\item Solve \eqref{eq:E1} over valid $\mathbb Z_2$ cocycles.  In the chosen gauge,
      $\sigma_{ij}=h_{1ij}$ for $i,j>1$.
\item Reconstruct $K_{ii}=d_i$ and $K_{ij}=\sigma_{ij}m_{ij}$, then verify $K^2=K$ exactly.
\item Apply Lemma~\ref{lem:zss}.  If it passes, the recovered gauge orbit is the unique one.
\end{enumerate}
The method is exponential in its most direct form; Theorem~\ref{thm:main} is an
identifiability result, not a polynomial-time learning theorem.  A useful local certificate
can nevertheless decode one edge.  If known legs give an anchored sum $a$ and the sum of
the absolute weights of all unknown legs is $u$, then $|a|>u$ fixes the sign of the full
row.  This dominance test certifies a warm-started solution but need not start from a tree
gauge by itself.

\section{Relation to principal-minor and boundary data}

General symmetric-kernel recovery from principal minors is governed by signed cycles and
principal-minor assignment \citep{HoltzSturmfels2007,GriffinTsatsomeros2006,
UrschelEtAl2017,Brunel2018,BrunelUrschel2024}.  Those results explain why two-point data do
not determine a generic symmetric kernel.  Theorem~\ref{thm:main} is compatible with them:
it restricts the matrix to the projection variety, where idempotency supplies the missing
cycle equations.  Our proof does not establish stability under noisy data, sparse-support
recovery, or the Hermitian phase analogue.

A subsequent related preprint by \citet{Villacis2026} uses corresponding principal
minors of all orders for complex orthogonal projections to determine the coordinatewise
absolute-value image of their ranges. The present theorem assumes only order-one and
order-two data and restricts recovery to generic real projections up to diagonal-sign
gauge. The different data requirements, domains, and recovery targets should be retained
in any novelty comparison; this observation does not establish priority.

There is a limited gravitational application.  In a finite real determinantal model, the
$n$-point correlation function is the determinant of the $n\times n$ kernel restriction.
Laplace-transformed one- and two-boundary data therefore determine $d_i$ and $m_{ij}$ in a
finite discretisation.  If the completion is also known to be a real projection, the
theorem says that, generically, these data fix the kernel up to a gauge that leaves all
principal minors unchanged.  Higher correlators then follow.  This is only a dictionary
for a finite real pure determinantal completion.  It does not prove that a gravitational
path integral has such a completion, nor does it settle Airy-kernel continuum limits,
complex phases, interactions, or ultraviolet completion.

\section{Limitations and conclusion}

The exceptional set is explicit but not classified: it is contained in the union of the
entry-zero and subset-sum-zero analytic loci.  Failure of the sufficient certificate does
not imply a rival projection.  Hermitian kernels replace $\mathbb Z_2$ holonomies by
$U(1)$ phases, so the discrete subset-sum proof does not transfer.  Continuum kernels
require analytic control of nodal domains and integral, rather than finite, wedge sums.

Within its declared scope, the result is complete.  The diagonal and entrywise moduli of a
generic interior real projection determine its signs up to diagonal gauge.  The proof
reduces purity to a holonomy-linear system, supplies a finite certificate, and establishes
that the certificate fails only on a measure-zero set.

\section*{AI-assisted workflow and unresolved attribution}
The repository operator states:
\begin{quote}
I genuinely do not know how authorship, intellectual credit, or ownership should be attributed for these results.
\end{quote}
The reported workflow was to ask AI systems to survey the literature,
autonomously select open problems, and attempt to solve them. The operator
reports that the earlier work was led mainly by Fable 5, whereas results from
July 2026 onward were predominantly produced by GPT models. DeepSeek and
multiple GPT models also participated in cross-review, checking, and
manuscript preparation and organization.

This is the operator's reported chronological history, not a complete audited
allocation of individual contributions. An earlier result is not reattributed
merely because it was revised or packaged after July. Human author names and
affiliations are intentionally blank. These results have not undergone human
expert review; automated checks and AI reviews are further AI review, not human
peer review. Historical source dates, present packaging dates, and demonstrable
public-disclosure dates are recorded separately. No authorship, ownership,
journal acceptance, or publication-priority determination is made here.

\bibliographystyle{plainnat}
\bibliography{references}
\end{document}
